% !TEX root =  ../mainPPDP.tex
%!TEX spellcheck = en_GB

The design of communication protocols for concurrent and distributed systems is a central concern in programming language research. 
Modern systems —- from microservices to IoT networks -— rely on structured sequences of message exchanges between multiple parties. 
In the family of behavioural types~\cite{DBLP:conf/concur/VasconcelosH93}, multiparty session types (MPST), as introduced by Honda et al.~\cite{HondaYC08,10.1145/2873052,DBLP:conf/icdcit/YoshidaG20}, have emerged as a powerful methodology to specify and verify such structured interactions. 
An MPST defines a global type describing the communication protocol as a whole, which can be projected into local types for each participant.% (as opposed to binary session types \cite{10.1007/3-540-57208-2_35,10.1007/3-540-58184-7_118,10.1007/BFb0053567} that only describe binary exchanges).

This type-based approach enables static verification of communication safety (no message type mismatches) and liveness properties (no deadlocks) by ensuring each party's code conformes to its local type. 
In essence, MPSTs allow protocol designers to treat communication patterns as first-class abstractions, yielding a rigorous framework for building correct distributed applications. 

In this work, we focus on the realisability problem, i.e., whether the local types projected from a global type can be recomposed into a concrete system that conforms to the intended global protocol\footnote{In classical MPST papers, this problem is also referred to as the implementability or projectability problem. We preferred to stick to realisability as a reference to the work by Alur et al. \cite{DBLP:journals/tcs/AlurEY05}.}. In classical MPST theory, protocols are often formulated under the sender-driven choice assumption: only one designated participant (the sender) decides between alternative message flow branches for all parties. 
While this assumption simplifies theoretical analysis, it proves too restrictive for many real-world protocols. 
Modern distributed applications typically involve multiple independent agents that can initiate interactions concurrently, requiring more flexible interaction patterns than a single centralised decision maker.
To illustrate the limitation of sender-driven choice, consider a simplified publish/subscribe scenario inspired by the MQTT (Message Queuing Telemetry Transport) protocol~\cite{mqtt311}.  Suppose we  have two client
\begin{wrapfigure}{r}{0.25\textwidth}
  \begin{center}
%\begin{figure}
    \begin{tikzpicture}[>=stealth, node distance=2cm, auto]
        % État initial
        \node[state, initial, initial text={}, initial distance=3mm, minimum size=1.5em] (s0) {\tiny$s_0$};
    
        % États pour publication c1 -> c2
        \node[state, minimum size=1.5em] (s1) [above right=2cm and 1cm of s0] {\tiny$s_1$};
        \node[state, minimum size=1.5em] (s2) [right=2cm of s0] {\tiny$s_2$};
    
        % Boucle de publication de c1
        \draw[->, bend left=40] (s0) to node[above,sloped]{\tiny$\gtlabel{c1}{b}{pub(m)}$}(s1);
        \draw[->, bend right=20] (s1) -- node[above,sloped] {\tiny$\gtlabel{b}{c2}{pub(m)}$} (s0);

        % Boucle de publication de c2
        \draw[->, bend left=20] (s0) to node[above,sloped]{\tiny$\gtlabel{c2}{b}{pub(m)}$}(s2);
        \draw[->, bend right=20] (s2) -- node[below,sloped]{\tiny$\gtlabel{b}{c1}{pub(m)}$} (s0);
    \end{tikzpicture}
    \caption{\label{fig:MQTT}MQTT protocol with two clients}
     \end{center}
\end{wrapfigure}
%\end{figure}
  processes (publisher/subscribers) $c_1$ and $c_2$ connected to a broker $b$. 
Both clients can independently publish messages to the broker, which in turn forwards each message to the other client. % (see Figure~\ref{fig:MQTT}). 
After the initial setup (e.g., subscriptions to topics not described here), either $c_1$ or $c_2$ may send a message at any time, without a predefined order. 
This pattern cannot be  captured by a sender-driven global type, as the choice of which client sends a message next is not controlled by any of the  participants.  In a sender-driven MPST, one would have to arbitrarily fix one client as the decision maker for the communication branch, which fails to reflect the  nature of the interaction.
%This example motivates the need to go beyond sender-driven protocols and to develop a theory of implementability for more general multiparty communications.

% MQTT 2-clients protocol global type

\subsubsection*{Contributions}In this work, we approach the realisability question from a language design perspective. We consider two  fundamental communication models for implementing MPST protocols: the asynchronous peer-to-peer ($\ppmodel$) model and the synchronous one.
In the $\ppmodel$ model, processes communicate via point-to-point asynchronous message passing over FIFO channels, which is the standard mode for distributed systems.
In the synchronous model, communications are rendezvous-based (handshakes where send and receive happen simultaneously), representing a globally synchronised message-passing semantics. 
These two models are widely considered as canonical extremes in analysing distributed protocols. 
A large amount of our approach generalises to even more communication models, but a full generalisation is beyond the scope of this paper,
and we only reason on general communication models for purpose of defining a notion of realisability that is parametric in the communication model (see Definition~\ref{def:realisability}).

In this work, we also investigate for the first time the question of the complementability of global types, 
i.e., whether a global type can be complemented by another global type that describes all possible complementary behaviours.
This notion may seem at first sight rather theoretical and unrelated to realisability, but our work reveals deep relations between
these two notions.

Our main contribution is to show that realisable global types form a subclass of the complementable ones.
More in details, we establish the following results:

\begin{itemize}
\item Every global type realisable in the synchronous communication model is complementable (Theorem~\ref{thm:realisable-complementable}), and the complementation is
effective (based on projection and NFA complementation)
\item Realisability is decidable in PSPACE for complementable global types provided an explicit complement is also given as input,
both in the synchronous communication model (Theorem~\ref{thm:decidability-of-implementability-in-synch}) and in the peer-to-peer model 
(Theorem~\ref{thm:decidability-of-implementability-in-p2p}), the latter result being based on recent results on RSC systems~\cite{germerie-phd}.
\item On the way, we show that if a global type is realisable in the peer-to-peer model, then it is also realisable in the synchronous one (Theorem~\ref{thm:pp-realizability-implies-synch-realizability}), which is not obvious, as 
deadlock-freedom is not a property that is monotonic in the communication model (a more permissive communication model may both hide some deadlocks and trigger new ones)
\item We collect several interesting observations about the question of the complementability of global types: not all global types are complementable (Example~\ref{ex:non-complementable-gt}),
but global types with at most three participants are complementable (Corollary~\ref{coro:3-participants-complementable}), 
as well as global types with sender-driven choices (Theorem~\ref{thm:sender-driven-choice-complementation}). For the latter, we provide
a complementation procedure with a linear blowup in the size of the global type.
\end{itemize}


% The communication models mentioned before form a hierarchy from relaxed to more constrained communication models, either in terms
% of executions~\cite{DBLP:conf/fm/ChevrouH0Q19}, or in terms of message sequence 
% charts~\cite{DBLP:journals/pacmpl/GiustoFLL23}.
% Several interesting properties turns out to be monotonic in the communication model in the following sense: if a property is satisfied in a communication model, then it is preserved for more relaxed (resp. more constrained) ones. 
% This  is true, for instance, when checking for reachability of a configuration,
% for unspecified receptions (the first message in a queue cannot be received by its
% destination process), orphan messages (i.e., messages sent by a peer but never received), etc.
% Unfortunately, implementability is not monotonic.
% To better understand this, take the following  two examples. First, consider an example of a global type (in standard MPST syntax) that is implementable in the peer-to-peer model but not in the mailbox one:  $$\gt_1 = \gtlabel{p}{q}{\msg_1} . \gtlabel{r}{q}{\msg_2} . \mathsf{end}.$$
%  This global type is implementable in the peer-to-peer model, 
% because the projection $p\parallel q \parallel r= !\msg_1 \parallel ?\msg_1.?\msg_2\parallel !\msg_2$ is deadlock-free and session conformant. However, this global type is not implementable in the mailbox model, where all incoming messages of process $q$ are stored in a unique FIFO queue,
% because $\msg_2$ could be sent and hence stored in the queue before $\msg_1$  and cause a deadlock.
% Second, consider this example where the converse happens: 
% %%
% $$ 
% \begin{array}{lll}
% \gt_2 & = & \gtlabel{p}{q}{\msg_1} . \textcolor{red}{\gtlabel{p}{r}{\msg_2}} . \gtlabel{p}{q}{\msg_3} 
%         . \textcolor{blue}{\gtlabel{q}{r}{\msg_4}} . \mathsf{end} 
% \\
% & + & \gtlabel{p}{q}{\msg_1'} . \textcolor{blue}{\gtlabel{q}{r}{\msg_4}} . \gtlabel{q}{p}{\msg_3'} 
%         . \textcolor{red}{\gtlabel{p}{r}{\msg_2}} . \mathsf{end}  
% \end{array}
% $$
% %%
% This global type $\gt_2$ is depicted in Figure~\ref{fig:gt2}, it features two possible
% scenarios, starting either with $\msg_1$ or $\msg_1'$, (that  are depicted on the left). 
% It is not implementable in the peer-to-peer model, because the projection 
% $p\parallel q \parallel r$ also admits the interaction where  $p$ and $q$ follow the first scenario, but $r$ follows the second scenario, resulting in a non synchronisable interaction (also depicted in Figure~\ref{fig:gt2}). However, this global type is implementable in the mailbox model, because the messages sent to $r$ are stored in a unique FIFO queue, and their order in the queue, is sufficient for  $r$ to  not confuse the two scenarios.
% \begin{figure}
%     \begin{tikzpicture}[scale=0.7]
%         \begin{scope}
%             \newproc{0}{p}{-2.7};
%             \newproc{2}{q}{-2.7};
%             \newproc{4}{r}{-2.7};

%             \newmsgm{0}{2}{-0.3}{-0.3}{1}{0.5}{black};
%             \newmsgm{0}{4}{-1.0}{-1.0}{2}{0.25}{red};
%             \newmsgm{0}{2}{-1.7}{-1.7}{3}{0.5}{black};
%             \newmsgm{2}{4}{-2.4}{-2.4}{4}{0.5}{blue};
%         \end{scope}

%         \begin{scope}[xshift=7cm]
%             \newproc{0}{p}{-2.7};
%             \newproc{2}{q}{-2.7};
%             \newproc{4}{r}{-2.7};

%             \newmsgm{0}{2}{-0.3}{-0.3}{1'}{0.5}{black};
%             \newmsgm{2}{4}{-1.0}{-1.0}{4}{0.5}{blue};
%             \newmsgm{2}{0}{-1.7}{-1.7}{3'}{0.5}{black};
%             \newmsgm{0}{4}{-2.4}{-2.4}{2}{0.25}{red};
%         \end{scope}

%         \begin{scope}[xshift=14cm]
%             \newproc{0}{p}{-2.7};
%             \newproc{2}{q}{-2.7};
%             \newproc{4}{r}{-2.7};

%             \newmsgm{0}{2}{-0.3}{-0.3}{1}{0.5}{black};
%             \newmsgm{0}{4}{-1.0}{-2.4}{2}{0.25}{red};
%             \draw[>=latex,->] (0,-1.45) -- node [below, pos = 0.25] {$\msg_{3}$} (2,-1.45);
%             \newmsgm{2}{4}{-1.95}{-1.95}{4}{0.5}{blue};
%         \end{scope}

%     \end{tikzpicture}
%     \caption{The two interactions specified by $\gt_2$, and another unspecified interaction that happens in the projected system in the p2p semantics}\label{fig:gt2}
% \end{figure}
% %%


% These two examples show that "implentability is not monotonic",
% and suggest that every new communication model should prompt the development of a new ad hoc theory of implementable global types, with its own syntactic conditions.



%\cinzia{add precision to contributions and fix Outline}
\subsubsection*{Outline.} The paper is organised as follows: Section~\ref{sec:preliminaries} introduces the necessary background on executions, Message Sequence Charts, communication models and communicating finite state machines. 
MPST and realisability are introduced in Section~\ref{sec:global_types}. Then Section \ref{sec:complementability} discusses the notions of effectively complementable and give examples of known global types that follow in this class. Finally Section~\ref{sec:decidability}  shows  our results on the reduction of realisability to the synchronous communication model based on the notion of effectively complementable global types.  
Section \ref{sec:concl} concludes with some final remarks and discusses related works.
%An appendix containing proofs and additional material has been included for the reviewer’s convenience.



